% Meeting 16
	\pagebreak
	\section*{Jan. 13, 2025}
	\addcontentsline{toc}{section}{Meeting 16}

	\textbf{C.A.V.E Attendees:} Abdul, Tabish, Berk, Andrew, Nicholas\\
	\textbf{Notetaker:} Nicholas

	\subsection*{Meeting with Dr. Giamou to plot Rev0}
		\subsubsection*{Agenda}
			\begin{itemize}
				\item Rev0 deliverable
			\end{itemize}

		\subsubsection*{Minutes}
			\begin{itemize}
				\item Supports unicorn or front position of the lidar on the helmet
				\item Need to avoid it sensing the helmet, unicorn avoids a lot of that
				\item Has the stencils for the PCBs- I(Nicholas) can use that tomorrow
				\item Mounting tabs- we can 3d print direct mount to the flatter pieces with tape
				\item Have some docker fixes, should be able to do 1-2 commands to get image working with Ros2
				\item Test ROS sync of both TOFs, can assume caver may not move between fast enough captures (25hz?)
				\item Gaussian processes- geostatistical analysis thing- assume signal is gaussianly related through time? smoothes clock jitter
				\item Have makerspace and solar car resources for space
				\item If we need a new pi, can order on the capstone dime
				\item ICP rundown- use both TOF, treat as single super sensor, run at 5hz, little overlap is fine, may be useful for calibration
				\item Batch offline state estimation- run algo at home
				\item IMU is ego motion frame, 'centre' gives accel and angular velocity
				\item motion/odometry from IMU, TOF, use laser to map onto space
				\item IMU as estimate for ICP, short time period integrated accel- open loop method
				\item Or, use as complementary filter (double accel integration has high error) assume accel is dominated by gravity, use angular velocity for orientation, roll and pitch without yaw
				\item May be existing ROS packages to do this, GTSAM is a python/cpp package for it 
				\item Solves opt problem of going between time steps from ICP Ti frames and complementary filters from IMU
				\item Perception in the Dark paper from 2020 uses the first method we considered
				\item Use up our budget first, and then Giamou can spot us for mmore if needed
				\item May meet on next Tuesday, depends where our work is, can meet at the same time every Tuesday if needed 
				\item 
			\end{itemize}

		\subsubsection*{Action Items}
			\begin{itemize}
				\item Use stencils to create some more boards
				\item Fix names of meeting minutes to make my life easier
				\item Read over paper 
				\item Test out mounting and capture
				\item 
			\end{itemize}


